\section{Introduction}
\label{sec:intro}

Who descends from whom is the central question of developmental biology, microbiology and cancer research.
Imaging answers it by tracking: segment cells in every frame of a time-lapse movie and link them across consecutive frames~\cite{ulman2017objective,maska2023ctc}.
Linking is easy only when frames are dense relative to cell motion and division.
Dense imaging, however, costs light dose, storage and instrument time, and many experiments cannot provide it at all: long experiments limited by phototoxicity, high-throughput screens that revisit each field only occasionally, or archives in which most frames were never kept.

We ask how much lineage can be recovered from the sparsest possible record: \emph{two snapshots} of the same population separated by $H$ minutes.
At $H{=}33$ minutes in a dense bacterial colony roughly every second cell has divided; at $H{=}132$ minutes each ancestor has up to eight descendants that have been pushed away by the growth of hundreds of neighbours (\cref{fig:teaser}).
The task is neither tracking (there is no smooth trajectory to follow) nor registration (the two point sets differ in size and topology).
We show that it is nonetheless learnable, and that it never needs annotated lineage: the learning signal can come from simulation, or from unannotated time-lapses labelled under a binary-fission constraint.

\paragraph{Family registration.}
A descendant family moves as a unit: its members are siblings and cousins packed together, displaced by the same pressure field.
\ours therefore predicts, for every ancestor, the centre of its descendant family at time $b$ and refines that centre through $T$ rounds of local cross-attention, in the spirit of iterative flow refinement~\cite{teed2020raft}.
Targets are scored against candidate ancestors by their offset from the refined centres; freezing the centres in an otherwise identical network costs 0.10 accuracy at 33 and 66 minutes.
A single network is trained across horizons from 8 to 240 minutes.

\paragraph{Learning lineage without lineage labels.}
Lineage annotations are scarce, and a model trained on one experiment transfers poorly to others (\cref{tab:external}).
A model trained only on \rodsim, an agent-based simulator of growing rod-shaped colonies, transfers well: it beats ten positions-only geometric, optimal-transport, registration and simulation baselines on all 27 bacterial sequences of four real datasets it was never trained on, by 3 to 10 points of accuracy over the best of them (three-seed mean); a plain learned matcher trained on the same simulations does as well, so the transfer is the simulator's.
A pre-registered blind test (protocol and hypotheses fixed in a file before the data were downloaded) confirms this on 28 unseen colonies grown in another medium, against the same baselines; earlier blind tests delimit its scope: it assumes a contiguous population expanding by its own growth, and fails on fields of many separate microcolonies and on steady-state chambers.
The remedy comes from biology.
Between consecutive frames, binary fission reduces lineage to an assignment in which each cell either persists or is replaced by daughters on either side of its centre along its axis; chaining such links through a dense \emph{unannotated} time-lapse yields pseudo-lineage at horizons up to four hours, 0.88--0.98 correct on the One-in-a-Million (OIAM) movies.
Trained on pseudo-lineage of three unannotated OIAM movies, \ours comes within 0.02--0.05 of the annotated model over all windows; decoded, it beats Trackastra~\cite{gallusser2024trackastra} fine-tuned on annotated lineage at 33 and 66 minutes and matches it at 132.
Bootstrapped self-labelling from the other half of a dataset's movies, fixed in advance, improves the simulation-trained model on all seven sources of our blind tests (pre-registered on the last three; on the four earlier ones, simpler Hungarian links did even better).
One dense, unannotated movie per organism could thus replace dense imaging later.

\paragraph{Genealogy in a single image.}
An oracle analysis shows that, given the true partition of targets into families, assignment accuracy at $H{=}132$ rises from 0.52 (our best model) to 0.78: errors pick a neighbour of the true ancestor, so the bottleneck is grouping, not coarse motion.
Grouping is visible in one frame: a network that predicts the time since two neighbours' lineages split, from a single phase-contrast image, detects same-family pairs with 0.97/0.91/0.70 average precision at $H{=}33/66/132$.
Coupling these affinities with \ours in a Potts mean-field decoder improves every base model at every gap.
The same network can be trained on \rodsim alone once the simulator randomises how cells divide, giving a pipeline with no real lineage label, though it then relies mostly on geometry.

\paragraph{Contributions.}
(1)~We formalise \emph{two-snapshot ancestry} and assemble a benchmark spanning six bacterial and human-cell datasets and three horizons, with cell-level accuracy and family-level biological readouts.
(2)~We propose \ours, an iterative family-registration network that resolves crowded colonies, where a learned matcher of the same capacity trails it by 0.10.
(3)~We show that \ours needs no annotated lineage: simulated colonies, or fission-constrained pseudo-lineage from unannotated time-lapses, suffice, including in pre-registered blind tests.
(4)~We predict genealogy from a single image and use it in a partition-aware decoder.
